\chapter{Annexes}
\label{chap:annexes}

\section{Raccourcis clavier}

\begin{tabularx}{\linewidth}{lX}
  \textbf{Raccourci} & \textbf{Action} \\
  \hline
  \texttt{Ctrl+1} & Accueil \\
  \texttt{Ctrl+2} & Plantations \\
  \texttt{Ctrl+3} & Cultures \\
  \texttt{Ctrl+4} & Lieux \\
  \texttt{Ctrl+5} & Calendrier \\
  \texttt{Ctrl+6} & Strates \\
  \texttt{Ctrl+7} & Paramètres \\
  \texttt{F1} ou \texttt{Ctrl+8} & Ouvrir ce manuel \\
  \texttt{Gauche / Droite} & Mois précédent/suivant (sur le Calendrier) \\
  \texttt{Tab} & Champ suivant dans un formulaire \\
  \texttt{Shift+Tab} & Champ précédent \\
  \texttt{Entrée} & Activer le bouton focusé \\
\end{tabularx}

\section{Glossaire}

\begin{description}
  \item[Annuelle] Plantation dont le cycle de vie tient sur une saison (semis $\to$ récolte $\to$ retrait).
  \item[Pluriannuelle] Plantation qui reste en place plusieurs années consécutives.
  \item[Variété] Déclinaison spécifique d'une culture (ex. « Coeur-de-bœuf » pour la tomate).
  \item[DTT] \emph{Days To Transplant} — nombre de jours entre le semis serre et le repiquage au champ.
  \item[DTM] \emph{Days To Maturity} — nombre de jours entre le repiquage et le début de récolte.
  \item[DOY] \emph{Day Of Year} — numéro de jour dans l'année (1 = 1\textsuperscript{er} janvier, 365 = 31 décembre).
  \item[Strate] Couche de végétation caractérisée par sa hauteur.
  \item[Lieu] Emplacement physique de culture (parcelle, planche, serre, verger, rang).
\end{description}

\section{Où obtenir de l'aide}

\begin{itemize}
  \item Site du projet : \url{https://guycorbaz.github.io/pomone}
  \item Issues GitHub : \url{https://github.com/guycorbaz/pomone/issues}
  \item Documentation technique : \texttt{pomone.pdf} dans la même release \texttt{docs-latest}.
\end{itemize}
